\section{Two half-rate bounds for directed recovery}
\label{sec:dc-section}

The distinction between a coordinate and its two incidence ports is essential
here.  A directed recovery relation on a finite set $V$, with $N=|V|$, assigns
a set $\mathcal N(q)\subseteq V\setminus\{q\}$ to each $q\in V$.
Its \emph{paired incidence graph} $E$ has ports $L_p,R_p$, $p\in V$, and
an undirected edge $L_pR_q$ precisely when $p\in\mathcal N(q)$.
Distances below are in this graph; disconnected ports have distance infinity.
Write
\begin{equation}
  b_E(r)=\bigl|\{p\in V:\operatorname{dist}_E(L_p,R_p)\le r\}\bigr|.
  \label{eq:dc-bad-pairs}
\end{equation}
A sequence $(E_n)$ on $2N_n$ ports, with $N_n\to\infty$, has
\emph{paired-distance separation} if $b_{E_n}(r)=o(N_n)$ for every fixed
integer $r\ge1$.  This concerns prescribed pairs of ports, rather than
the girth or connected-component sizes of $E_n$.

We prove a half-rate bound under either of two additional hypotheses.
Uniform hyperfiniteness suffices for every finite-alphabet code.
Without hyperfiniteness, bounded complex linear recovery coefficients
and a fixed block width suffice.  The latter extends the graph hypothesis
of Theorem~A of~\cite{BrendeCurved} from a locally bipartite underlying
graph to paired-distance separation.

\subsection{All codes on hyperfinite incidence graphs}

A storage code over a finite alphabet $\mathcal A$ of size at least two is a set
$\mathcal C\subseteq\mathcal A^V$ such that two codewords agreeing on
$\mathcal N(q)$ agree at $q$, for every $q\in V$.
Equivalently, each coordinate has a recovery function on the neighbour
assignments that occur in the code.  The code need not be linear, and
no distribution or extension property is assumed.

\begin{lemma}[Separated components]
\formal{IncidenceComponents.lean}{storageCode_card_le_of_separated_components}
\label{lem:dc-components}
If $L_p$ and $R_p$ lie in different components of $E$ for every $p\in V$,
then every storage code satisfies
\[
  |\mathcal C|\le |\mathcal A|^{\lfloor N/2\rfloor}.
\]
\end{lemma}
\begin{outline}
Let $\alpha(p)$ and $\beta(p)$ be the components containing $L_p$ and $R_p$.
Choose an order of the components, and record the coordinates $p$ for which
$\alpha(p)$ precedes $\beta(p)$.  This record determines a codeword.
Indeed, process the $\alpha$-components in increasing order.  A coordinate
in the current component was either recorded, or its $\beta$-component is
earlier.  In the latter case every recovery neighbour $u\in\mathcal N(p)$
satisfies $\alpha(u)=\beta(p)$, so its value is already determined.
The equality case $\alpha(p)=\beta(p)$ is excluded by hypothesis.
Compare this order with its reverse. Every coordinate is recorded in exactly
one of the two orders, so one record has at most $\lfloor N/2\rfloor$
coordinates. Its projection is injective on $\mathcal C$.
\end{outline}

\begin{theorem}[Finite separator bound]
\formal{IncidenceSeparator.lean}{storageCode_separator_log_bound}
\label{thm:dc-separator}
Let $Z$ be a set of incidence ports such that each component of $E-Z$
has at most $K$ vertices, where $K\ge1$.  Every nonempty storage code satisfies
\begin{equation}
 \log_{|\mathcal A|}|\mathcal C|
 \le \frac{N+|Z|+b_E(K)}2.
 \label{eq:dc-separator-bound}
\end{equation}
No degree bound is required.
\end{theorem}
\begin{outline}
Let $U$ contain the coordinates with at least one port in $Z$, and also
the coordinates counted by $b_E(K)$.  Thus $|U|\le |Z|+b_E(K)$.
Delete both ports of every coordinate in $U$, obtaining the induced
relation on $Q=V\setminus U$.  Its components still have at most $K$
ports.  If a surviving pair were connected, a shortest path in its
component would have length at most $K-1$, contradicting the choice of $U$.

Fix an assignment to the coordinates in $U$.  The restrictions to $Q$ of
codewords with this assignment form a storage code for the induced relation:
deleted recovery arguments are fixed constants.  \cref{lem:dc-components}
bounds each such fibre by $|\mathcal A|^{|Q|/2}$.  There are at most
$|\mathcal A|^{|U|}$ fibres, giving
$|\mathcal C|\le |\mathcal A|^{|U|+|Q|/2}$ and the claimed inequality.
This conditions on coordinate values; it does not discard side-information
edges without paying for their coordinates.
\end{outline}

\begin{definition}
\label{def:dc-hyperfinite-definition}
\formal{Hyperfinite.lean}{UniformlyHyperfinite}
A sequence $(E_n)$ is \emph{uniformly hyperfinite} if, for each $\eta>0$,
there is an integer $K_\eta$ such that, for every sufficiently large $n$,
one can delete at most $2\eta N_n$ incidence ports and leave components
of size at most $K_\eta$.
\end{definition}

\begin{corollary}[Hyperfinite half-rate bound]
\formal{Hyperfinite.lean}{hyperfinite_rate_limsup_le_half}
\label{cor:dc-hyperfinite}
Suppose $(E_n)$ is uniformly hyperfinite and has paired-distance separation.
For any finite alphabets $\mathcal A_n$ of size at least two, and any
nonempty storage codes $\mathcal C_n$ on these relations,
\[
 \limsup_{n\to\infty}
 \frac{\log_{|\mathcal A_n|}|\mathcal C_n|}{N_n}\le\frac12.
\]
\end{corollary}
\begin{outline}
For fixed $\eta$, \cref{thm:dc-separator} gives a limiting upper bound
$1/2+\eta$, because $K_\eta$ is fixed before $n$ tends to infinity.
Then let $\eta$ decrease to zero.  The finite estimate is uniform in
the alphabet, so it also permits the displayed varying alphabets.
\end{outline}

\subsection{Bounded complex linear recovery}

Let $t\ge1$.  A recovery matrix $A$ has $t\times t$ blocks $A_{qp}$,
with $A_{qq}=0$ and $A_{qp}=0$ unless $L_pR_q\in E$.
Its actual support graph has $L_pR_q$ precisely when $A_{qp}\ne0$;
we may use a larger declared relation $E$.
A linear storage space with these recovery rules is contained in
$\ker(I-A)$.  No symmetry or adjoint relation between $A_{qp}$ and
$A_{pq}$ is imposed.

\begin{theorem}[Directed complex half-rank bound]
\notformalized
\label{thm:dc-complex}
Fix integers $\Delta,t\ge1$ and a real number $B\ge1$.
Let $(E_n)$ have maximum incidence degree at most $\Delta$ and paired-port
separation.  Let $A_n$ be complex recovery matrices of block width $t$
on $E_n$, with every scalar entry of modulus at most $B$.
Then
\begin{equation}
 \liminf_{n\to\infty}
 \frac{\operatorname{rank}_{\mathbb C}(I-A_n)}{tN_n}\ge\frac12.
 \label{eq:dc-complex-rank}
\end{equation}
Consequently every associated complex or real linear storage space has
dimension at most $(1/2+o(1))tN_n$.
\end{theorem}

We give the analytic details because two points matter: a component
marking must preserve unimodularity, and finite nullity passes to the
limit in only one direction.
We use the marked-network framework of Aldous and Lyons~\cite[Sections~2
and~5]{AldousLyons}.  Uniformly rooted finite graphs of bounded degree,
with marks in fixed compact spaces, have locally convergent subsequences;
their limits satisfy the mass-transport identity.  For a unimodular law
$\mu$, bounded measurable operator families equivariant under network
isomorphisms, and independent of the distinguished root, form a von Neumann
algebra.  In its $t\times t$ matrix amplification the normalized trace is
\begin{equation}
 \tau(T)=\mathbb E_\mu\frac1t\operatorname{tr}(T_{oo}),
 \qquad \tau(I)=1.
 \label{eq:dc-trace}
\end{equation}
It is tracial and positive.  These are the same analytic inputs used in
the proof of Theorem~A in~\cite{BrendeCurved}.

\begin{informalproof}[Proof of \cref{thm:dc-complex}]
Form the coordinate graph $G_n$ by joining $p$ and $q$ if either directed
incidence is present.  Its degree is at most $2\Delta$.
Record the two directed incidence flags as discrete marks, and the
corresponding bounded coefficient blocks as compact marks.  In particular,
retain an incidence flag when its coefficient converges to zero.
Choose a subsequence realizing the limsup of the normalized nullities,
then a locally convergent marked subsequence with unimodular limit $\mu$.
The root is uniform among coordinates, and the rooted limit retains the
root's connected coordinate component.  Write $(G,A,o)$ for that limit.
The row and column absolute sums of $A$ are at most $\Delta tB$, so
\[
 \|A\|\le\Delta tB,\qquad \|I-A\|\le c:=1+\Delta tB.
\]

\emph{Separation in the limit.}
For fixed $r,s$, the event that a coordinate within $G_n$-distance $s$
of the root has paired-port distance at most $r$ has probability at most
\[
 \left(\sum_{j=0}^s(2\Delta)^j\right)\frac{b_{E_n}(r)}{N_n}.
\]
It is determined by finitely many discrete incidence flags and therefore
is a clopen local event.  Its limiting probability is zero.  Taking the
countable union over $r,s$ shows that, almost surely, \emph{every}
coordinate $p$ of $G$ has $L_p,R_p$ in different components of the limit
port graph $E$.

\emph{An invariant phase extension.}
Conditionally on this countable graph, give each component of $E$ an
independent uniform phase on the complex unit circle.  This law can be
specified without choosing an invariant enumeration: for a finite list
of ports, give independent phases to its distinct component classes,
and equal phases to ports in the same class.  The cylinder laws are
consistent.  Connectivity of two ports is a countable union of finite-path
events, so their component partition is Borel.  The cylinder integrals,
and then all bounded Borel integrals by a monotone-class argument, define
a Borel probability kernel $\kappa_G$.  Its law is unchanged by any
isomorphism or change of root.

The resulting phase-marked extension $\widehat\mu$ is unimodular.
Indeed, for any nonnegative doubly rooted test function $f$, average
$f$ over $\kappa_G$ to obtain $\bar f$.  The same kernel is used with
either root distinguished.  Applying mass transport to $\bar f$ and
Tonelli's theorem gives
\[
 \mathbb E_{\widehat\mu}\sum_x f(G,o,x)
 =\mathbb E_\mu\sum_x\bar f(G,o,x)
 =\mathbb E_\mu\sum_x\bar f(G,x,o)
 =\mathbb E_{\widehat\mu}\sum_x f(G,x,o).
\]
This construction is a probability-space extension, not an assertion
that component phases are a factor of independent vertex labels.

Let $D_L,D_R$ multiply the $t$ entries at coordinate $p$ by the phase
of the component of $L_p,R_p$, respectively.  They are measurable,
equivariant diagonal unitaries.  An edge $L_pR_q$ equates the phase
at column $p$ with that at row $q$, and hence
\begin{equation}
 AD_L=D_RA.
 \label{eq:dc-intertwining}
\end{equation}
The phases of each matched pair are independent and nonatomic.
By countability they differ at every coordinate simultaneously, almost
surely.  Thus $D_L-D_R$ is injective; no bounded inverse is asserted.

\emph{The tracial half bound.}
Let $K=\ker(I-A)$ and let $P_K$ be its orthogonal projection.
If $y\in K$ and $D_Ly\in K$, then
$D_Ly=AD_Ly=D_RAy=D_Ry$, so $y=0$.
Consequently $T=(I-P_K)D_LP_K$ has kernel $K^\perp$, and its range
closure lies in $K^\perp$.
All these operators lie in the amplified equivariant algebra.
Polar decomposition there gives a partial isometry $U$ with
$U^*U=P_K$ and $UU^*\le I-P_K$.  One can obtain $U$ as the strong
limit of $T(T^*T+\varepsilon I)^{-1/2}$ as $\varepsilon\downarrow0$;
functional calculus and strong closure keep this limit in the algebra.
Therefore
\begin{equation}
 \tau(P_K)=\tau(U^*U)=\tau(UU^*)\le1-\tau(P_K),
 \qquad \tau(P_K)\le\frac12.
 \label{eq:dc-half-trace}
\end{equation}
This uses the range closure, not an assumption that the range is closed.
The expected root trace of $P_K$ is unchanged on forgetting the phases.

\emph{Return to finite matrices.}
Put $M_n=I-A_n$, $H_n=M_n^*M_n$, and $H=(I-A)^*(I-A)$.
Their normalized empirical spectral measures $\nu_n$ have common support
in $[0,c^2]$.  For each integer $j\ge0$, the moment
$(tN_n)^{-1}\operatorname{tr}H_n^j$ is the expected normalized root-block
trace of a bounded continuous local function; radius $2j$ suffices.
Local convergence gives moment convergence to the averaged root spectral
measure $\nu$ of $H$.  Polynomial approximation on $[0,c^2]$ gives weak
convergence.  The closed-set inequality for the singleton $\{0\}$ is
\begin{equation}
 \limsup_n\frac{\dim\ker M_n}{tN_n}
 =\limsup_n\nu_n(\{0\})
 \le\nu(\{0\})=\tau(P_K)\le\frac12.
 \label{eq:dc-nullity-semicontinuity}
\end{equation}
Here $\ker(M^*M)=\ker M$ holds for finite matrices and bounded operators.
This proves~\eqref{eq:dc-complex-rank}.  Real matrices have the same rank
after complexification.
\end{informalproof}

\begin{corollary}[Uniform finite form]
\notformalized
\label{cor:dc-uniform}
For fixed $\Delta,t,B$ and every $\varepsilon>0$ there exist an integer
$r\ge1$, a number $\delta>0$, and $N_0$ such that
\[
 b_E(r)\le\delta N,\quad N\ge N_0
 \quad\Longrightarrow\quad
 \operatorname{rank}_{\mathbb C}(I-A)\ge(1/2-\varepsilon)tN
\]
for every relation and matrix with the bounds of \cref{thm:dc-complex}.
\end{corollary}
\begin{informalproof}
Otherwise choose a failure with $r=j$, $\delta=1/j$, and $N\ge j$
for each positive integer $j$.  These failures have paired-distance separation,
contradicting \cref{thm:dc-complex}.
\end{informalproof}

\begin{corollary}[Complex linear codes on actual cone relations]
\notformalized
\label{cor:dc-native}
Fix the tape count, linear-time constant, full cone depth, block width~$t$,
and a uniform bound on the recovery coefficients.  Every sequence of
complex linear storage spaces on the corresponding native cone relations,
admitting linear recovery maps with the stated block width and coefficient
bound, has dimension at most $(1/2+o(1))tL$ as $d\to\infty$.
\end{corollary}
\begin{informalproof}
Remove the $O(L/\log d)$ self-demand coordinates supplied by
\cref{cor:nb-self-demands}.  For any fixed $0<\delta<1$, also remove
coordinates whose left degree exceeds a fixed $M\ge B_D/\delta$, where
$B_D=\sum_{h=0}^D(k+1)^h$ is the cone bound in
\eqref{eq:nb-cone-bound}, using
\cref{cor:nb-trim}.  Write $H$ for their union and $Q$ for the
retained set.  Then $|H|\le\delta L+o(L)$, and the induced relation is
loop-free, has bounded incidence degree, and satisfies the paired-distance
condition by \cref{thm:nb-main}.

If $\mathcal U$ is the original linear space, its subspace with zero
coordinates on $H$ has codimension at most $t|H|$.  On this subspace the
original bounded linear recovery maps restrict to maps on $Q$ with the
same coefficient bound and block width.  \cref{thm:dc-complex} bounds
its dimension by $(1/2+o(1))t|Q|$.  Therefore
\[
 \dim\mathcal U\le t|H|+(1/2+o(1))t|Q|
 \le(1/2+\delta/2+o(1))tL.
\]
Take the size limit first and then let $\delta$ tend to zero.  This is a
complex-linear statement about the cone relations; it does not make the
machine's binary genie a complex linear map.
\end{informalproof}

\begin{remark}[Scope of the graph extension]
\label{rem:dc-comparison}
For a symmetric recovery relation, $E$ is its bipartite double cover.
A path between matched ports projects to an odd closed walk, and an odd
cycle gives such a path at each of its vertices.  With bounded degree,
paired-distance separation is consequently equivalent to having only $o(N_n)$
vertices on odd cycles of every fixed bounded length.
For directed relations the hypotheses differ: in disjoint directed
seven-cycles with one recovery neighbour at each vertex, $E$ is a union
of isolated edges and every matched distance is infinite, whereas the
underlying coordinate graphs have odd cycles at positive density.
Thus replacing directed recovery by its undirected closure loses information.
\end{remark}

\subsection{What the complex bound says over finite fields}

\begin{corollary}[Necessary torsion in bounded lifts]
\notformalized
\label{cor:dc-torsion}
Fix a prime $p$, integers $\Delta,t\ge1$, and $\varepsilon>0$.
Let $(E_n)$ satisfy the graph hypotheses of \cref{thm:dc-complex}.
Suppose matrices $M_n$ over $\mathbb F_p$ have diagonal blocks $I_t$,
off-diagonal support in $E_n$, and
$\dim\ker M_n\ge(1/2+\varepsilon)tN_n$.
For any integer lifts $\widetilde M_n$ with the same support restrictions,
literal diagonal blocks $I_t$, and uniformly bounded entries, at least
$(\varepsilon-o(1))tN_n$ of their \emph{nonzero} Smith invariant factors
are divisible by $p$.
\end{corollary}
\begin{informalproof}
Set $\rho_n=\operatorname{rank}_{\mathbb Q}\widetilde M_n$.
\cref{thm:dc-complex} gives $\rho_n\ge(1/2-o(1))tN_n$.
The unimodular transformations in Smith normal form remain invertible
modulo $p$.  Thus
\[
 \#\{\hbox{nonzero invariant factors divisible by }p\}
 =\rho_n-\operatorname{rank}_{\mathbb F_p}M_n
 \ge(\varepsilon-o(1))tN_n.
\]
Equivalently, the subgroup annihilated by $p$ in the integer cokernel
has at least this dimension over $\mathbb F_p$.  Zero invariant factors
give free abelian summands and are not counted as torsion.
\end{informalproof}

This conclusion permits such torsion; it is not a finite-field half-rate
theorem.  A fixed extension field can be represented over its prime field
by enlarging the fixed block width.  The analytic argument does not cover
growing block widths, unbounded complex coefficients, or nonlinear codes
on nonhyperfinite relations.  \cref{sec:mc-section} gives a separate
finite-field certificate on one existential expanding family, uniform
in the storage-code block width.  Neither result identifies all native
cone relations with that family or settles the unrestricted coding problem.
